% ============================================================
%  RESUMEN / ABSTRACT
% ============================================================

Este trabajo presenta CAIQ (\textit{Career Alignment and Insight Qualifier}), un
sistema de recomendación profesional orientado a personas que se encuentran en proceso
de transición hacia perfiles del sector de datos. El sistema toma como punto de
partida el currículum vitae del candidato, extrae sus competencias, las contrasta con
los requisitos reales de ofertas de empleo del rol objetivo y cuantifica la distancia
existente. A partir de ese análisis, genera recomendaciones personalizadas de
formación ---cursos y másteres--- y de empleo, ordenadas por su contribución al cierre
de la brecha identificada.

El sistema se articula en cuatro componentes. Un \textit{pipeline} de \textit{scraping}
incremental mantiene actualizado un corpus de ofertas de empleo, cursos y programas de
máster procedentes de portales especializados. Un módulo ETL normaliza y enriquece esos
datos sobre una taxonomía controlada de competencias del sector. Un módulo de modelado
del perfil extrae las competencias del CV, infiere el nivel de \textit{seniority} y
calcula el \textit{skill gap} respecto al rol objetivo, estructurado en bandas de
prioridad. Finalmente, un motor de recomendación híbrido combina coincidencia léxica y
similitud semántica para ordenar los recursos más relevantes en cada caso.

El trabajo cubre todo el ciclo: desde la recopilación y preprocesamiento de los datos
hasta el despliegue de una aplicación web de acceso público, pasando por el diseño del
motor de recomendación, su comparación entre distintas variantes de \textit{scoring} y
su validación sobre casos de uso representativos. El resultado es un sistema funcional,
validado y desplegado, que ofrece al profesional en transición una referencia objetiva
sobre dónde está y qué necesita aprender.

\vspace{0.4em}
\noindent\textbf{Palabras clave:} sistemas de recomendación, \textit{skill gap},
mercado laboral de datos, extracción de competencias, NLP, orientación profesional,
\textit{scraping}, ETL.

% ── Abstract (English) ───────────────────────────────────────────────────────
\vspace{2em}
\begin{otherlanguage*}{english}
\section*{Abstract}
\markboth{Abstract}{Abstract}
\addcontentsline{toc}{section}{Abstract}

This work presents CAIQ (\textit{Career Alignment and Insight Qualifier}), a
professional recommendation system designed for people transitioning towards data roles.
The system takes the candidate's CV as its starting point, extracts their competencies,
contrasts them against the real requirements of job postings for the target role and
quantifies the existing gap. From this analysis, it generates personalised
recommendations for training ---courses and master's programmes--- and employment,
ranked by their contribution to closing the identified gap.

The system is built around four components. An incremental scraping pipeline keeps a
corpus of job postings, online courses and master's programmes up to date from
specialised portals. An ETL module normalises and enriches these data against a
controlled competency taxonomy for the data sector. A profile modelling module extracts
competencies from the CV, infers the candidate's seniority level and computes the skill
gap with respect to the target role, structured in priority bands. A hybrid
recommendation engine then combines lexical matching and semantic similarity to rank
the most relevant resources for each candidate.

The work covers the full cycle: from data collection and preprocessing through
recommendation engine design, multi-variant scoring comparison and validation on
representative use cases, to the deployment of a publicly accessible web application.
The outcome is a functional, validated and deployed system that gives the transitioning
professional an objective reference for where they stand and what they need to learn.

\vspace{0.4em}
\noindent\textbf{Keywords:} recommender systems, skill gap, data job market, skill
extraction, NLP, career guidance, web scraping, ETL.
\end{otherlanguage*}
